% =====================================================================
%  Цель и задание лабораторной работы.
% =====================================================================

\textbf{Цель:} освоить на практике основные принципы реализации
информационно-поисковых систем и методы оценки качества их работы.

\vspace{1em}

\textbf{Задание (вариант 22)}: разработать информационно-поисковую систему
со следующей комбинацией признаков варианта:
\begin{itemize}
  \item \textbf{Углублённый элемент} -- модуль индексирования документов
  (построение поискового образа документа: нормализация текста, расчёт
  TF, IDF, весов терминов и вектора документа);
  \item \textbf{Сфера применения} -- сеть Интернет (документы поступают
  через собственного поискового робота, а не вводятся вручную);
  \item \textbf{Стратегия поиска} -- векторная модель (TF-IDF и косинусная
  мера между вектором документа и вектором запроса);
  \item \textbf{Язык} -- английский.
\end{itemize}

\vspace{1em}

Согласно методическим указаниям, система должна:
\begin{enumerate}
  \item принимать на входе множество естественно-языковых текстов, по
  которым осуществляется поиск, и автоматически выделять ключевые слова
  документов в соответствии с формулой веса термина;
  \item позволять пользователю формулировать запрос на естественном языке
  и строить по нему поисковый образ запроса (ПОЗ);
  \item выдавать список документов, ранжированных по релевантности
  запросу в соответствии с векторной моделью поиска, с активной ссылкой
  на документ и списком слов запроса, найденных в документе;
  \item на основании метрик, применяемых для оценки качества
  информационно-поисковых систем (методика ROMIP), программно вычислять
  и отображать в виде таблиц и графиков оценку качества работы системы;
  \item предоставлять простой и доступный пользовательский интерфейс со
  справочными подсказками.
\end{enumerate}
